% Part: methods
% Chapter: induction
% Section: strong-induction

\documentclass[../../../include/open-logic-section]{subfiles}

\begin{document}

\olfileid{mth}{ind}{str}

\olsection{પ્રબળ આગમન}

ઉપર ચર્ચેલા આગમનના સિદ્ધાંતમાં આપણે $P(0)$ સાબિત કરીએ છીએ
અને જો $P(k)$ હોય તો $P(k+1)$ પણ હોય છે એ સાબિત કરીએ છીએ.
બીજા ભાગમાં $P(k)$ સાચું છે એમ ધારીને આ ધારણા વડે
$P(k+1)$ સાબિત કરીએ છીએ. અલબત્ત, તેની સમકક્ષ રીતમાં
$P(k-1)$ ધારીને $P(k)$ સાબિત કરી શકીએ. મુખ્ય વાત એ છે કે
કોઈ પણ સંખ્યાથી તેના અનુગામી સુધીનું તારણ કાઢી શકીએ:
કોઈ સંખ્યા માટે દાવો તેના પૂર્વગામી માટે સાચો છે એમ
ધારીને તે સંખ્યા માટે સાબિત કરી શકીએ.

આગમનના સિદ્ધાંતનો એક પ્રકાર એવો છે જેમાં $k$ના પૂર્વગામી
$k-1$ માટે જ દાવો સાચો છે એમ નથી ધારતા;~$k$થી નાની
બધી સંખ્યાઓ માટે તે સાચો છે એમ ધારીએ છીએ અને તે
ધારણાથી~$k$ માટે દાવો સ્થાપીએ છીએ. આથી પણ બધા~$n \in
\Nat$ માટે $P(n)$ મળે છે. કારણ કે એક વાર $P(0)$
સ્થાપીએ, ત્યારે~$1$થી નાની બધી સંખ્યાઓ માટે $P$ સાચો
છે એવું સ્થાપ્યું હોય છે. અને જો બધા $l<k$ માટે $P(l)$
હોય તો $P(k)$ થાય છે એમ જાણતા હોઈએ, તો ખાસ કરીને
$k=1$ માટે પણ આ જાણીએ છીએ. તેથી $P(1)$ તારવી
શકીએ. હવે $P(0)$ અને $P(1)$ સાબિત થયા, એટલે બધા
$l<2$ માટે $P(l)$ સાબિત થયું. બધા $l<2$ માટે $P(l)$
હોય તો $P(2)$ થાય છે એવું શરતી વિધાન પણ હોવાથી
$P(2)$ તારવી શકીએ, અને આ રીતે આગળ વધીએ.

ખરેખર, જો "બધા $k$ માટે, બધા $l<k$ માટે $P(l)$
હોય તો $P(k)$" એવું સામાન્ય શરતી વિધાન સ્થાપી શકીએ,
તો $P(0)$ અલગથી સ્થાપવાની જરૂર નથી, કારણ કે તે
શરતી વિધાનમાંથી જ નીકળે છે. યાદ રાખો કે "બધા
$l<k$ માટે $P(l)$" એવો સામાન્ય દાવો $l<k$ હોય
એવી કોઈ સંખ્યા ન હોય ત્યારે પણ સાચો છે. આ
ખાલી પરિમાણીકરણનો કિસ્સો છે: "$A$ પ્રકારની
બધી વસ્તુઓ $B$ પ્રકારની છે" એ જો $A$ પ્રકારની
કોઈ વસ્તુ ન હોય તો સાચું છે; કોઈ $x$ જો~$!A(x)$
સંતોષતું ન હોય, તો $\lforall[x][(!A(x) \lif !B(x))]$
સાચું છે. અહીં ઔપચારિક સ્વરૂપ
"$\lforall[l][(l < k \lif P(l))]$" થશે; અને જો
$l < k$ હોય એવી કોઈ સંખ્યા ન હોય તો તે સાચું છે.
હવે $k=0$ હોય ત્યારે બરાબર એવું જ થાય છે:
કોઈ $l<0$ નથી; તેથી "બધા $l<0$ માટે $P(l)$"
સાચું છે, $P$ ગમે તે હોય.\sourcecorrection{OLMD-164}{સ્થિર અંગ્રેજી મૂળમાં
  અહીં $P(l)$ને બદલે $P(0)$ લખાયું છે. ખાલી ક્ષેત્રે બંને
  સાર્વત્રિક વિધાનો સાચાં છે, પરંતુ આગમનના પૂર્વપદમાં
  દરેક $l<k$ માટે $P(l)$ જ જરૂરી છે.} તેથી "બધા
$l<k$ માટે $P(l)$ હોય તો $P(k)$" એવી સાબિતી
આપમેળે~$P(0)$ પણ સ્થાપે છે.

જો~$k$ માટે દાવો સ્થાપવા માત્ર $k-1$ માટેનો દાવો
પૂરતો ન હોય અને એક કે વધુ $l<k$ માટેનો દાવો
ધારવો પડે, તો આ પ્રકારનું આગમન ઉપયોગી છે.

\end{document}
